% ============================================================
% CHAPTER: LIMITATIONS AND FUTURE WORK
% Content moved here from the Discussion (limitations) and the
% Conclusion (future work) so that the scope of the claims and the
% research agenda are visible in the table of contents rather than
% buried as subsections. Labels are preserved so cross-references
% elsewhere continue to resolve.
% ============================================================

The claims in this thesis are bounded, and stating those bounds precisely is
essential to not overstating them. This section groups the limitations into those
of \emph{scope} (what was varied) and those of \emph{method} (how the motifs were
measured), and then maps them onto a prioritised research agenda.

\subsection{Limitations}
\label{sec:discussion:limitations}

\paragraph{Scope: what we did and did not vary.}
The experimental grid varies communication topology, memory window, task
difficulty, and benchmark, and the seeded experiment additionally varies the
opening vote composition causally. It does \emph{not} vary the base model (a
single model family throughout), the number of agents ($N = 4$), agent
heterogeneity (all agents share priors and competence), or the interaction
framing (cooperative, synchronous, two-phase). Every quantitative claim should be
read as conditional on these fixed choices. This matters most for the results that
already differ between our two benchmarks: if a single model on two datasets
produces opposite signs for the value of dominance and of stylistic diversity,
then the effect of changing the model, or introducing capability asymmetry, is
likely to be at least as large. The heterogeneity axis is especially consequential
because the critical literature identifies model heterogeneity, not debate per se,
as the reliable source of gains \citep{liu2025stopovervaluing, chen2024reconcile},
and capability asymmetry is what drives the accuracy-decreasing herding reported by
\citet{wynn2025talk}. Our homogeneous design isolates the dynamics of a uniform
collective by construction, but it cannot speak to heterogeneous ones.
Recent evidence tempers this expectation in one direction. In a controlled
comparison of mixed against single-model teams within the 7-8B class,
heterogeneous compositions underperformed their best homogeneous member in six of
eight condition-task pairs, and modal sycophancy in mixed teams remained
comparable to homogeneous teams \citep{bertalanic2026costofconsensus}. Model
diversity therefore does not appear to be sufficient on its own to dissolve the
dynamics we characterise, which strengthens rather than weakens the case for
first characterising the homogeneous baseline. Whether the same holds across
capability tiers rather than within one remains untested.

A related scope boundary concerns the memory window. We tested $W \in \{1, 2, 5\}$,
a range that keeps the history short relative to the round cap. At sufficiently
large $W$, every agent can see the complete history of every other agent's messages
regardless of topology: a spoke in a star can reconstruct the full exchange through
accumulated hub summaries, and the structural difference between fully-connected and
star communication becomes negligible. This suggests that the topology effects we
report---faster convergence under FC, hub-driven limit cycles under star---could
diminish or vanish if $W$ were large enough to bridge the information gap between
topologies. Our results characterise the low-to-moderate memory regime; whether
the topology effect survives large memory windows is an open question.

\paragraph{Decoding stochasticity.}
Every result in this thesis is conditioned on a single decoding temperature
($T = 0.4$). Temperature controls the entropy of the output distribution: lower
temperatures concentrate probability mass on the model's modal completion, higher
temperatures spread it across alternatives. In the context of debate dynamics
this is a first-order parameter: at high temperature agents start with more
diverse positions and are harder to lock into a fixed point; at low temperature
they converge faster but may become overconfident. A recent physics-inspired
account models the biased consensus problem as a Potts spin system and shows
analytically that the system undergoes a phase transition as temperature increases
from a consensus phase to a disordered phase \citep{okawa2026biasedconsensus}.
Empirically, ablating temperature from $T = 0.4$ to $T = 0.7$ in a comparable
protocol reduced the modal sycophancy rate by $6.2$ percentage points
\citep{bertalanic2026costofconsensus}, suggesting that our findings may be
conservative: the self-reinforcement and convergence rates we report are likely
higher than they would be at typical deployment temperatures. Replicating the
main analyses at higher temperatures is a natural sensitivity check we did not
perform.

\paragraph{Task-selection bias from Tier-3 screening.}
All quantitative results are conditioned on Tier-3 tasks, those where the outcome
varies across repetitions under the screening configuration.
By construction, this excludes tasks that are trivially easy (unanimous at $t=0$),
nearly deterministic (the group always converges to the same answer), or too hard
for any signal to emerge.
The prevalence rates reported here, how often each motif appears, are therefore
rates \emph{within the contested regime}, not across the full task distribution.
Extrapolating them to the broader population of questions, or to operating
conditions where most tasks fall outside Tier~3, would overstate motif frequency.

\paragraph{The causal evidence is GPQA-only.}
The seeded-initialisation experiment, our only intervention and the strongest
support for the amplifier claim, was run exclusively on GPQA
(\autoref{sec:results:seeded}). This is a real restriction on the causal claim,
and in one respect the less favourable case is the one left untested: on
HiddenBench, where agents begin systematically wrong, amplification of the opening
distribution should if anything be more consequential (and more often harmful),
so the mechanism is plausibly at least as strong there. But that remains an
expectation, not a measurement. The causal statement should therefore be read as
established for GPQA and as a well-supported but observational hypothesis for
HiddenBench; a seeded HiddenBench arm is the most valuable single extension of this
work.

\paragraph{Adversarial agents and deployment risk.}
The experiments in this thesis use cooperative, homogeneous agents: every agent is
instructed to reason honestly and update in good faith.
A devil's advocate probe (\autoref{sec:results:seeded}) suggests this is a
consequential assumption.
A single agent instructed to advocate persistently for a wrong answer reduces
system accuracy from $63.7\,\%$ to $10.0\,\%$ even when two of four agents start
correct, falling below the $26.2\,\%$ obtained with a natural $1/4$-correct
minority and no adversary.
This is a deployment risk: a consistently biased agent---whether from a
misconfigured prompt, a compromised model, or adversarial injection---exploits the
same front-loaded update dynamics that allow a lone correct minority to occasionally
flip the group.
Because the protocol has no mechanism for distinguishing persistent conviction from
persistent misinformation, the amplifier can be hijacked asymmetrically: truth-seeding
produces a $13.0\times$ lift relative to no ground truth; adversarial persistence
erases that lift and then some.
Quantifying this vulnerability across group sizes and topologies, and evaluating
structural defences---outlier detection via vote entropy, trust discounting for
non-updating agents---is an important direction that this thesis leaves open.

\paragraph{Effective topology analysis is FC-only and threshold-based.}
The effective influence topology is extracted only from FC runs; the star topology's
influence graph is structurally prescribed and the classification is not
meaningful there.
Within FC runs, the four-type taxonomy (flat, hub, two-hub, star) is defined by
threshold rules on the share of total influence captured by the top one or two
agents.
The star boundary ($\geq$75\,\% for a single agent) and the sample sizes for the
rare types, 2.3\,\% and 5.8\,\% star rates, mean that type-level estimates carry
substantial uncertainty.
The design-guideline claim (diverse init raises P(star)) is observational: we
measure that P(star) co-varies with initial vote diversity across naturally-occurring
runs but do not directly manipulate diversity.

\paragraph{Method: signatures, not certified mechanisms.}
All four \suitename{} detectors identify \emph{signatures}, not certified dynamical
objects. The self-reinforcement slope is a phenomenological positive trend, not a
proof of a positive-feedback loop; the RQA test flags periodic recurrence, not a
mathematically certified limit cycle; $N_{\mathrm{eff}}$ measures endpoint
coexistence, not a certified multistable attractor landscape; and the dominance
index measures conversion concentration, not certified causal influence. With the
single exception of the seeded-initialisation experiment, which is a genuine
intervention, every association reported here is observational. Establishing that,
say, hub identity \emph{causes} the outcome, rather than co-varying with a
question-specific expertise signal, would require perturbation experiments we did
not run.

\paragraph{Specific measurement caveats.}
Four narrower caveats bound specific results. (i) The self-reinforcement null may
be subject to survivorship bias: agents whose confidence drifts downward are more
likely to flip their vote and terminate a vote-stable segment, which can censor
negative slopes and inflate $p_{\mathrm{SR}}$ above the $0.5$ baseline; the
prevalence result is therefore conservative in direction but not bias-free in
magnitude. (ii) The diversity$\to$SR$\to$rounds mediation is partly mechanical,
since longer segments afford more rounds over which a slope can accumulate; we report
it as a decomposition, not as evidence that SR temporally precedes convergence.
(iii) The bistability $N_{\mathrm{eff}}$ measure can in principle exceed the option
count when split (non-unanimous) endpoints are admitted; because we restrict
attractors to unanimous endpoints (a $2.1\,\%$ loss of repetitions),
$N_{\mathrm{eff}}$ is bounded by the option count by construction, and we verified
$N_{\mathrm{eff}} \leq M$ in all $420$ cells.
(v) The dominance influence graph is built from vote adoptions during the
\updatephase{} phase, which is precisely where the protocol instructs agents to
update after reading peers' \arguephase{} drafts.
Conversion edges therefore partly reflect the structural fact that agents are
\emph{supposed} to be influenced in this phase; a high dominance score on a
hub-to-spoke edge may indicate genuine persuasive leadership or simply that the
protocol routes all information through the hub.
We treat dominance as a concentration \emph{signature} and never as a certified
causal influence for this reason.

\paragraph{Statistical multiplicity.}
This is an exploratory study reporting many tests across twelve configuration
cells and several motifs. Inference is conducted at the task level (the unit of
analysis, $n \approx 35$; \autoref{sec:method}), and we apply corrections within
families of tests where appropriate (Benjamini--Hochberg false-discovery-rate
control across configuration cells and outcome variables, permutation and
surrogate nulls per detector), but we do not impose a single study-wide
correction, and some false positives are expected across the full battery. We have
tried to lean on effect sizes and on consistency across cells rather than on
isolated $p$-values, and the central claims (universality of SR, the amplifier
effect, the topology/accuracy dissociation) rest on effects that are large and
repeated across many cells rather than on any single borderline test.


\subsection{Future Work}
\label{sec:conclusion:future}

The limitations above map directly onto a prioritised research agenda.

\begin{enumerate}
  \item \textbf{From signatures to causes.} With the exception of the seeded
        experiment, our findings are observational. Targeted perturbation
        experiments, seeding basins to certify the bistability landscape,
        perturbing the hub to test whether dominance is causal or an expertise
        proxy, and an asynchronous protocol variant to isolate synchrony-induced
        cycles, would upgrade the remaining associations toward causal claims.

  \item \textbf{Effective topology under heterogeneity and adversarial conditions.}
        The effective topology analysis is confined to FC runs with homogeneous
        agents. Whether effective star formation rates change under model
        heterogeneity, capability asymmetry, or partial non-cooperation is
        unknown; the devil's advocate result suggests the formation mechanism is
        exploitable, so characterising it under adversarial conditions is a
        high-priority extension. A targeted diverse-initialisation intervention
        that deliberately seeds each agent with a distinct starting answer would
        also test the P(star) lever causally rather than observationally.

  \item \textbf{Heterogeneity and scale.} The most consequential untested axes are
        model heterogeneity and group size. The critical literature identifies
        heterogeneity as the reliable source of debate gains
        \citep{liu2025stopovervaluing, chen2024reconcile} and capability asymmetry
        as a driver of harmful herding \citep{wynn2025talk}; testing whether the
        amplifier characterisation survives heterogeneous agents and larger $N$ is
        the natural next step, and one our homogeneous design cannot address.

  \item \textbf{Correcting the self-reinforcement null.} The survivorship bias in
        the SR prevalence null (downward-drifting confidence censored by vote
        flips) should be addressed with a matched null that models the flip hazard,
        to establish how much of the above-baseline $p_{\mathrm{SR}}$ is genuine.

  \item \textbf{Toward the diagnostic suite.} Finally, the \suitename{} instruments
        in \autoref{tab:design:instruments} could be packaged into an online monitor
        and the Tier-2 recommendations validated across models, group sizes, and
        task domains beyond the two studied here. This is the step that would turn
        the groundwork of this thesis into the practical tool it points toward.

  \item \textbf{A generative model of debate dynamics.} The motif taxonomy is
        descriptive. A forward model of how agents update---whether as Fictitious
        Play, gradient-based reinforcement, or a social-learning rule---would make
        testable predictions about which motifs arise under which conditions and
        would connect the empirical patterns to game-theoretic equilibrium concepts
        \citep{mehdizadeh2026topologymemory}. This is the missing link between the
        diagnostic layer and a prescriptive theory of multi-agent deliberation.

  \item \textbf{Removal-based influence attribution.} The dominance metric
        identifies concentration by tallying vote adoptions, a phenomenological
        measure. Complementing it with removal-based attribution---estimating each
        agent's marginal causal contribution to the final outcome by masking or
        ablating their messages---would give a sharper, interventional measure of
        leadership \citep{agentsthatmatter2026, cui2025introspecloo, infodynamics2026}.
        Applied to hub agents, this would test whether structural dominance translates
        to causal influence or whether dominance nodes are proxies for
        question-specific expertise.
\end{enumerate}
